% BEGIN REV09_RECTAS27_COORDENADAS
\subsection{Coordonnées réversibles des paires de droites}
\label{corpus9:xii:rectas-coordenadas}

Le tableau~\ref{p12-83-c20-s6:tab:h3-27-rectas-u029} fixe la bijection
\(\iota:H_3\to\mathscr L_{27}\), où \(\mathscr L_{27}\) est la
configuration des vingt-sept droites. Nous conservons la loi centrale
\eqref{p12-83-c20-s6:eq:heisenberg-3-nuevo} :
\[
 (a,b,z)(c,d,w)=(a+c,b+d,z+w+bc),\qquad
 (a,b,z)^{-1}=(-a,-b,-z+ab).
\]
La formule de l’inverse se vérifie en multipliant dans les deux ordres ;
les trois coordonnées obtenues sont nulles. Toutes les opérations de cette
partie s’effectuent dans \(\mathbb F_3\).

\begin{proposition}[Codage réversible des 729 mots]
\label{corpus9:xii:rectas-biyeccion}
Pour \(w=(\ell_1,h_1,\ell_2,h_2,\ell_3,h_3)\), soient
\(L(w)=(\ell_1,\ell_2,\ell_3)\) et
\(H(w)=(h_1,h_2,h_3)\). L’application
\begin{equation}
 \mathcal E:\mathbb F_3^6\longrightarrow
                 \mathscr L_{27}\times\mathscr L_{27},\qquad
 \mathcal E(w)=(\iota L(w),\iota H(w))
 \label{corpus9:xii:rectas-codificacion}
\end{equation}
est bijective. Dans les coordonnées relatives
\(L=(a,b,z)\), \(R=L^{-1}H=(u,v,t)\), sa reconstruction est

\begin{equation}
 H=(a+u,b+v,z+t+bu),\qquad
 w=(a,a+u,b,b+v,z,z+t+bu).
 \label{corpus9:xii:rectas-reconstruccion}
\end{equation}
\end{proposition}
\begin{proof}
Étant donnée une paire ordonnée \((\ell,h)\) de droites, le tableau fixe les
uniques triplets \(L=\iota^{-1}(\ell)\) et
\(H=\iota^{-1}(h)\). Entrelacer leurs coordonnées récupère un unique
mot \(w\), et l’évaluation de
\eqref{corpus9:xii:rectas-codificacion} renvoie la paire initiale.
Réciproquement, désentrelacer puis entrelacer de nouveau un mot conserve
ses six coordonnées. Cela établit les deux compositions inverses.

La transformation \((L,H)\mapsto(L,L^{-1}H)\) est bijective, d’inverse
\((L,R)\mapsto(L,LR)\). La loi centrale donne
\(LR=(a+u,b+v,z+t+bu)\), ce qui prouve
\eqref{corpus9:xii:rectas-reconstruccion}. En particulier, le terme
\(bu\) conserve le cocycle dans la troisième coordonnée. Si on l’omet,
la reconstruction change précisément de \(bu\) ; par exemple,
\(L=(0,1,0)\), \(R=(1,0,0)\) produisent
\(H=(1,1,1)\), tandis que la somme composante par composante produit
\((1,1,0)\).
\end{proof}

La partition du théorème~\ref{p12-83-c20-s6:thm:particion-hensel-schlafli-nueva}
est conservée sous cette bijection : l’élément relatif \(R\) distingue
la diagonale, l’incidence et le gauchissement, de cardinaux \(27\), \(270\)
et \(432\). Les couples ordonnés constituent le domaine complet du
système de coordonnées. La condition d’adjacence est une condition supplémentaire sur
les couples ou leurs suites. Le prolongement aux histoires complètes est
construit dans le théorème~\ref{corpus9:xii:rectas-homeomorfismo} au moyen
du lecteur de blocs du TPK.
% END REV09_RECTAS27_COORDENADAS
